\section{Availability}
\label{sec:availability}

\paragraph{Weights.}
The Bobcat 1 adapter evaluated in this report is not distributed. Its successors are released
on Hugging Face under Apache-2.0, each with its SHA-256 sums, release manifest, the identifier
list the compiler needs and a model card that reports its own evaluation: Bobcat 1.1, a new
rank-16 adapter on the same Qwen3.8-27B revision
(\url{https://huggingface.co/sanghwa-na/bobcat-1.1}; ready-to-serve NVFP4 weights at
\url{https://huggingface.co/sanghwa-na/bobcat-1.1-nvfp4}), and Bobcat Flash 1.1, a distilled
adapter on Gemma 4 26B-A4B-it (\url{https://huggingface.co/sanghwa-na/bobcat-flash-1.1};
merged weights at \url{https://huggingface.co/sanghwa-na/bobcat-flash-1.1-merged}). The
compiler, server, training and evaluation code are at \url{https://github.com/foxl-ai/bobcat}
under Apache-2.0. The derived evaluation data are not redistributed.

\paragraph{Internal API.}
The hosted API at \url{https://api.foxl.ai/bobcat} runs since 2026-09-25 and is restricted
to internal (administrator) accounts; it is not offered to the public. The paragraphs below
describe that deployment.

\paragraph{What the internal deployment serves.}
It serves the 27B model as the merged FP8 artifact of Section~\ref{sec:engine} on vLLM, on
one NVIDIA L40S 48GB (an H100 was chosen, but none was available; the L40S limits concurrent
sequences to 128). On development this artifact scores 93.1\%
task macro and agrees with the evaluated (unmerged) model on 98.3\% of answers
(Section~\ref{sec:engine_quality}); the sealed final result of 94.6\% was measured on the
unmerged path and is not a measurement of the served artifact.

\paragraph{Wire compatibility.}
The API uses the paths and shapes of TypeSafe's published HTTP contract: \texttt{POST
/v1/systemone} for decisions, \texttt{GET /v1/models} for the model list, and \texttt{GET
/health}. Existing TypeSafe SDK code can call Bobcat by setting the SDK's base URL to
\url{https://api.foxl.ai/bobcat} and using a foxl API key issued to an internal account.
We checked this with the unmodified \texttt{typesafe-sdk} 0.7.1 for Python against our API
server running with a deterministic stand-in engine and no weights: \texttt{system\_one}
with Noul, Choice and Score questions returned typed results, \texttt{models.list()}
returned the model list, and an invalid request (a Choice with no criteria) was rejected with
HTTP 422, which the SDK surfaced as \texttt{TypeSafeUnprocessableEntityError}. This checks the
wire format only, not decision quality. For compatibility the server accepts TypeSafe's model
aliases in the request, but the \texttt{model} field of every reply names the Bobcat version
that actually answered.

\paragraph{Operational behaviour.}
The origin server accepts only requests forwarded by an authenticating edge, which also
applies the rate limits. The server accepts up to 128
questions per request and 255 candidates per Choice, and refuses a question whose compiled
sequence exceeds 16,384 tokens with a JSON 422 instead of truncating it. Request bodies above
2\,MiB are refused (HTTP 413), engine failures return a fixed JSON 500, and overload returns
529 with a retry hint. Replies carry the number of processed tokens and the confidence profile
in response headers. The origin host is patched and rebooted weekly by an account-level
patching policy; for those few minutes the edge returns 529.

\paragraph{What callers should know.}
A typed reply can be the wrong valid option, and \texttt{confidence} is not a probability of
being right (Section~\ref{sec:confidence}). Arithmetic, date comparison, string matching,
permissions and side effects belong in the caller's code. For tasks with untrusted input we
recommend stating in the instructions that instructions inside the state are not evidence
(Section~\ref{sec:injection}); we have not measured whether this reduces the risk measured
there.

\paragraph{Provenance.}
The backbone, Qwen3.8-27B, is released by its authors under Apache-2.0 \citep{qwen38_27b}. The
evaluation uses KLUE data \citep{park2021klue} under CC BY-SA 4.0; rows derived from it are not
redistributed. Appendix~\ref{app:compute} lists the hashes that identify the evaluated
artifact.
